\section{Conclusion and future work}

This work designed and implemented an end-to-end computer vision pipeline aimed at the tactical digitisation of football matches. The resulting system is a working tool that combines several architectures to turn a video signal into information about the match. Along the way, the effectiveness of modern detection and tracking models was validated: architectures such as YOLO and association algorithms such as ByteTrack are able to keep player identities even under partial occlusion and dense movement.

A significant contribution of this work was the team segmentation module based on visual \textit{embeddings}. Using pre-trained vision models (DINOv2) together with dimensionality reduction (UMAP) and clustering (K-Means) proved to be a better strategy than traditional colour histogram methods. This approach separated the teams semantically with high accuracy and without any match-specific or league-specific retraining. The patch-based (slicing) inference strategy implemented for ball detection likewise gave impressive results when detecting and tracking small, fast-moving objects, even when they occupy a minimal fraction of the frame.

Evaluating the keypoint detection model on new video sequences, however, exposed critical generalisation challenges that had not shown up during controlled experimentation. The test videos did not belong to the same dataset used to train the keypoint models. This visual discrepancy between the training images and the new videos showed that the model was unable to generalise the geometric features of the pitch correctly to different leagues, stadiums or cameras. Together with the noise introduced by compression and camera motion, this caused fluctuations in the homography matrix that had not appeared during evaluation on the original test set.

As future work, we propose hardening the system by attacking the observed instability directly. A priority line of improvement is to analyse all the frames in order to assess the consistency of the homography across them and correct it in between, that is, to look at the behaviour over the whole video rather than at each frame and its neighbours only. It would also be interesting to add action detection models to the pipeline, so as to gain a better understanding of what is happening at each moment of the match.
